
\begin{table}[!h]
\centering
\caption{Descriptive Statistics}
\label{tab:1}
\resizebox{\ifdim\width>\linewidth\linewidth\else\width\fi}{!}{
\begin{tabular}[t]{llllllr}
\toprule
Variable & Group 1 & Group 2 & Group 3 & Group 4 & Total & p-value\\
\midrule
age & 22.24 (9.90) & 21.40 (8.44) & 21.82 (9.36) & 20.94 (6.62) & 21.59 (8.61) & 0.476\\
\addlinespace
male & 0.51 (0.50) & 0.46 (0.50) & 0.55 (0.50) & 0.41 (0.50) & 0.48 (0.50) & 0.476\\
\addlinespace
studied economics & 0.15 (0.36) & 0.15 (0.36) & 0.12 (0.33) & 0.10 (0.30) & 0.13 (0.33) & 0.355\\
\addlinespace
proxy household income & 10528.91 (4952.51) & 9790.42 (4860.53) & 9253.72 (5183.12) & 10077.93 (5293.96) & 9926.28 (5057.82) & 0.529\\
\addlinespace
financial literacy & 2.56 (0.69) & 2.44 (0.78) & 2.61 (0.62) & 2.63 (0.66) & 2.56 (0.69) & 0.319\\
\addlinespace
cb knowledge & 0.24 (0.56) & 0.38 (0.75) & 0.28 (0.49) & 0.27 (0.51) & 0.30 (0.59) & 0.938\\
\addlinespace
elicited 12 month ahead expectations & 6.05 (3.40) & 5.53 (3.67) & 5.90 (3.27) & 5.65 (3.73) & 5.77 (3.52) & 0.685\\
\addlinespace
prior trust & 2.66 (0.95) & 2.55 (0.91) & 2.84 (1.06) & 2.81 (1.01) & 2.71 (0.98) & 0.189\\
\addlinespace
\bottomrule
\end{tabular}}
\begin{minipage}{\textwidth}
    {\fontsize{8}{8}\selectfont\textit{Notes:} The table reports the mean and standard deviation (in parenthesis) for the demographics as well as the prior beliefs about the economy of the subjects in each treatment group. Household Income is the mean household income of the subjects' self reported neighborhood of residence. Trust in the Central Bank refers to trust in the BCB to care about the economic well-being of all Brazilians. The final column displays the p-values from a one-way Analysis of Variance (ANOVA) test, which compares the means across groups to verify whether they are significantly different.}
\end{minipage}
\end{table}

